%------------------------------------------------------------------------------%

\usepackage{calculo_varias_variaveis-1}

\title{Vetores e Planos}

%------------------------------------------------------------------------------%
\begin{document}

\addtitle

%------------------------------------------------------------------------------%
\section{Vetores}

%------------------------------------------------------------------------------%
\begin{frame}{Vetores em $R^3$}
  Segmento de reta orientado

  \pause
  Dois vetores são iguais se tem o mesmo comprimento e direção

  \pause
  Componentes
  \[
    v
    = \left< \nu_1, \nu_2, \nu_3 \right>
    = \left( \nu_1, \nu_2, \nu_3 \right)
  \]

  \pause
  Magnitude ou comprimento
  \[
    \abs{v} = \sqrt{\nu_1^2 + \nu_2^2 + \nu_3^2}
  \]

  \pause
  Vetor unitário
  \[
    \abs{v} = 1
  \]
\end{frame}

%------------------------------------------------------------------------------%
\begin{frame}{Produto Escalar}
  Produto escalar entre
  \(u=\left<u_1, u_2, u_3\right>\) e
  \(v=\left<v_1, v_2, v_3\right>\)
  \[
    u\cdot v = u_1v_1 + u_2v_2 + u_3v_3
  \]

  \pause
  Ângulo entre dois vetores
  \[
    \cos(\theta) = \frac{u\cdot v}{\abs{u}\abs{v}}
  \]

  \pause
  Vetores ortogonais
  \[
    u\cdot v = 0
  \]
  \end{frame}

%------------------------------------------------------------------------------%
\section{Retas e Planos}

%------------------------------------------------------------------------------%
\begin{frame}{Retas}
  Equação vetorial para uma reta
  \[
    r(t) = r_0 + tv
    \qquad
    t\in\R
  \]
  Reta que passa pelo ponto \(r_0=\left(x_0, y_0, z_0\right)\)
  na direção do vetor \(v=\left<v_1, v_2, v_3\right>\)

  \pause
  \vspace{\baselineskip}
  Equações paramétricas
  \begin{align*}
    x & = x_0 + tv_1 \\
    y & = y_0 + tv_2 \\
    z & = z_0 + tv_3 \\
  \end{align*}
\end{frame}

%------------------------------------------------------------------------------%
\begin{frame}{Planos}
  Definimos um plano por um ponto, $P_0=(x_0, y_0, z_0)$, \\
  \pause
  e uma inclinação (vetor perpendicular), $n=\left<A, B, C\right>$

  \pause
  \vspace{\baselineskip}
  $P=(x, y, z)$ pertence ao plano se
  \[
    n\cdot \overrightarrow{P_0P} = 0
  \]
  \pause
  \[
    \left<A, B, C\right>\cdot \left<x-x_0, y-y_0, z-z_0\right> = 0
  \]
  \pause
  \[
    A(x-x_0) + B(y-y_0) + C(z-z_0) = 0
  \]
  \pause
  \[
    Ax + By + Cz = Ax_0 + By_0 + Cz_0
  \]
\end{frame}

%------------------------------------------------------------------------------%
\section{Lista Mínima}

%------------------------------------------------------------------------------%
\begin{frame}{Lista Mínima}
  Cálculo Vol. 2 do Thomas 12\textsuperscript{a} ed.
  \begin{enumerate}
    \item Estudar o texto das Seções 12.2, 12.3, 12.5
  \end{enumerate}

  \vfill
  Atenção: A prova é baseada no livro, não nas apresentações
\end{frame}

%------------------------------------------------------------------------------%
\end{document}
%------------------------------------------------------------------------------%
